\label{chap:p3-83-45-completitud_finita}

\section{Extensión conservativa de la completitud finita mediante multiplicidades proyectivas}

Una torre HMT concreta está formada por conjuntos finitos, relaciones de
transición, aplicaciones de restricción y lectores. Todo ello puede
codificarse en ZF o ZFC. La colección de todas las torres constituye una
clase definible; no se necesita postular que HMT forme ya una teoría de
clases autónoma comparable con NBG.

\begin{proposition}[Codificación conservativa]
Cada objeto genealógico finito o numerable utilizado en este cuaderno puede
representarse mediante conjuntos y relaciones de ZFC. Por tanto, las
formalizaciones aquí construidas no demuestran una contradicción de ZFC.
\end{proposition}

\begin{proof}
Cada nivel \(X_m\) es un conjunto, cada restricción \(p_{n,m}\) es un
subconjunto de \(X_n\times X_m\), y una torre numerable es una función con
dominio \(\mathbb N\) que codifica esos datos. Los tipos, lectores y memorias
se añaden como relaciones o funciones. El límite inverso es un subconjunto
del producto \(\prod_mX_m\), que ZFC construye mediante Reemplazo, Unión,
Potencia y Separación.
\end{proof}

Lo que HMT discute no es la legitimidad de esa codificación, sino la práctica
de olvidar los mapas y conservar sólo una sombra. Si se aplican, por ejemplo,
el lector
\[
  L:\OmegaInf\longrightarrow\mathbb R
\]
o el perfil cardinal
\[
  \mathbb X\longmapsto(|X_m|)_{m\ge0},
\]
pueden colapsarse objetos no isomorfos. La proyección es válida en su
codominio; lo inválido sería atribuirle fidelidad sin prueba.

\begin{remark}[Condiciones de validez]
No se afirma que HMT satisfaga internamente todos los esquemas de ZF, IZF o
NBG. Separación, Reemplazo, Potencia plena y una jerarquía transfinita general
no han sido verificados conjuntamente dentro de una axiomática HMT autónoma.
\end{remark}

\section{Categoría de objetos genealógicos}

\begin{definition}
Un objeto genealógico es
\[
  \mathbb X=
  \bigl((X_m)_{m\ge0},(p_{n,m})_{n\ge m},
  \mathcal T,\mathcal O,\mathcal M\bigr),
\]
donde \(\mathcal T\) registra tipos, \(\mathcal O\) observables y
\(\mathcal M\) memorias. Un morfismo
\(f:\mathbb X\to\mathbb Y\) es una familia \(f_m:X_m\to Y_m\) que satisface
\[
  p^Y_{n,m}f_n=f_mp^X_{n,m}
\]
y preserva los datos señalados.
\end{definition}

La categoría así obtenida es una presentación concreta de objetos
pro-finitos enriquecidos. El funtor de secciones globales
\[
  \Gamma(\mathbb X)=\varprojlim_mX_m
\]
recupera las historias completas. La categoría pequeña de caminos de
\(\TPK\) admite además la envolvente de presheaves
\[
  \widehat{\mathcal P_{\TPK}}
  =
  \mathbf{Set}^{\mathcal P_{\TPK}^{op}}.
\]
Esta construcción proporciona lógica contextual y haces de datos, pero no
autoriza a declarar booleanidad, elección o suficientes puntos sin estudiar
el topos concreto.

\section{Fundación de prefijos y coinducción de secciones}

El grafo de prefijos está graduado por la profundidad:
\[
  \operatorname{rk}(x_m)=m.
\]
Aplicar el padre disminuye estrictamente el rango. No existen ciclos
parentales
\[
  x_m\to x_{m-1}\to\cdots\to x_m
\]
compatibles con esa graduación.

\begin{proposition}
La bien fundación del grafo de prefijos y la existencia coinductiva de una
sección infinita son compatibles.
\end{proposition}

\begin{proof}
La bien fundación concierne a la dirección de descenso por padres. La
coinducción concierne a una familia \((x_m)\) compatible con todos los
descensos. Ningún \(x_m\) pertenece a sí mismo ni reaparece como el mismo
nodo; la sección es una función sobre \(\mathbb N\), no un ciclo de
pertenencia.
\end{proof}

Esta distinción aclara el retorno \(9\to1\). La fase puede repetirse, pero
la coordenada de memoria aumenta. Un lazo en la proyección de fase no es un
ciclo en el grafo graduado de estados.

\section{Fibración de los ordinales por historias genealógicas compatibles}

Von Neumann representa cada ordinal como el conjunto de sus predecesores
\citep{vonneumann1923}. HMT puede elevar esa representación sin modificarla.

\begin{definition}[Elevación genealógica de un ordinal finito]
Sea \(\boldsymbol x=(x_m)\) una sección. Se define
\[
  \widehat n_{\boldsymbol x}
  =
  \{(k,x_k):k<n\}.
\]
\end{definition}

\begin{proposition}
La proyección sobre la primera coordenada satisface
\[
  \pi(\widehat n_{\boldsymbol x})=n,
\]
y la sucesión obedece
\[
  \widehat{n+1}_{\boldsymbol x}
  =
  \widehat n_{\boldsymbol x}\cup\{(n,x_n)\}.
\]
\end{proposition}

\begin{proof}
Ambas identidades se obtienen directamente de la definición. La información
añadida no altera el orden subyacente; distingue las rutas que realizan el
mismo ordinal.
\end{proof}

Para \(\omega\) se escribe
\[
  \widehat\omega_{\boldsymbol x}
  =
  \{(n,x_n):n\in\mathbb N\}.
\]
El soporte de cada historia sigue siendo numerable, aunque el espacio de
todas las historias pueda tener cardinal continuo cuando la ramificación
persiste.

\begin{remark}[Extensión y dominio abierto]
La extensión
\[
S_{\alpha+1}=\operatorname{Ext}_\alpha(S_\alpha),
\qquad
S_\lambda=\varprojlim_{\beta<\lambda}S_\beta
\]
es un programa transfinito bien tipado, no un resultado ya demostrado para
todo ordinal. Debe verificarse la coherencia de los límites, el tamaño de los
niveles y la existencia de extensiones.
\end{remark}

\section{Funtor de decategorización y cardinalidad de las secciones publicadas}

Dos torres pueden satisfacer
\[
  |X_m|=|Y_m|\qquad\text{para todo }m
\]
y no ser isomorfas: basta que sus mapas de restricción tengan patrones de
ramificación distintos. El perfil cardinal es, por tanto, un invariante
útil pero no fiel.

\begin{remark}[Interpretación estructural]
Contar cuántas casillas existen no dice qué casilla prolonga a cuál, qué
orientación se invierte, qué carry se transporta o qué estado retorna. El
cardinal mide el inventario; la genealogía mide la organización.
\end{remark}

\begin{remark}[Separación de resultados y alcance]
La codificación ZFC, la sección como límite y la distinción entre igualdad
observacional e identidad de estado son
\emph{resultados establecidos en las secciones precedentes}. La categoría unificada, el topos de
presheaves y el ordinal fibrado son
\emph{construcciones explícitas del presente desarrollo}. No se proclama una teoría autónoma completa
de clases o de ordinales transfinítos HMT.
\end{remark}


\section{Límite profinito del sistema de truncamientos genealógicos}

Sea
\[
  X=\varprojlim_mX_m
\]
un límite inverso de conjuntos finitos discretos con mapas sobreyectivos.
Dotamos \(X\) de la topología inicial de las proyecciones
\(\pi_m:X\to X_m\). Los cilindros
\[
  [A]_m=\pi_m^{-1}(A),
  \qquad A\subseteq X_m,
\]
son clopen y forman una base.

\begin{theorem}[Potencia cilíndrica]
\label{p3-83-c45-s2:thm:clopen-power}
\[
  \Clop(X)
  \cong
  \varinjlim_m\bigl(\mathcal P(X_m),p_{n,m}^{-1}\bigr).
\]
\end{theorem}

\begin{proof}
Todo cilindro es clopen. Recíprocamente, si \(C\subset X\) es clopen, para
cada \(x\in C\) existe un cilindro contenido en \(C\). La compacidad de \(C\)
da un subrecubrimiento finito. Elevando esos cilindros a una profundidad
común \(m\), se obtiene \(C=\pi_m^{-1}(A)\) para cierto
\(A\subseteq X_m\). Dos representantes en profundidades distintas definen el
mismo clopen exactamente cuando coinciden tras una elevación común.
\end{proof}

\begin{theorem}[Potencia cerrada]
\label{p3-83-c45-s2:thm:closed-power}
\[
  \Closed(X)
  \cong
  \varprojlim_m\bigl(\mathcal P(X_m),p_{n,m*}\bigr),
\]
donde \(p_*(A)=p(A)\).
\end{theorem}

\begin{proof}
Si \(F\subset X\) es cerrado, es compacto. Sus proyecciones
\(F_m=\pi_m(F)\) satisfacen \(p_{n,m}(F_n)=F_m\). Recíprocamente, una familia
compatible \((F_m)\) define el límite inverso
\[
  F=\{x\in X:\pi_m(x)\in F_m\text{ para todo }m\}.
\]
Es una intersección de cerrados y, por la compatibilidad sobreyectiva de las
restricciones \(F_n\to F_m\), sus proyecciones son exactamente los \(F_m\).
Ambas construcciones son inversas.
\end{proof}

\section{Obstrucción a la reconstrucción de la potencia genealógica a partir de reducciones finitas}

En un espacio de Cantor,
\[
 |\Clop(X)|=\aleph_0,
 \qquad
 |\Closed(X)|=2^{\aleph_0},
 \qquad
 |\mathcal P(X)|=2^{2^{\aleph_0}}.
\]
Los tres objetos no pueden identificarse.

Esta separación no presupone la hipótesis del continuo. Escribiendo
\(\mathfrak c=2^{\aleph_0}\),
\[
 |\Clop(X)|=\aleph_0,\qquad
 |\Closed(X)|=\mathfrak c,\qquad
 |\mathcal P(X)|=2^{\mathfrak c}.
\]
\(\mathsf{CH}\) sólo identifica \(\mathfrak c\) con \(\aleph_1\); no
colapsa ninguna de las tres capas. La independencia Gödel--Cohen cambia qué
cardinal puede ocupar \(\mathfrak c\), no la desigualdad de Cantor
\(\mathfrak c<2^{\mathfrak c}\).

\begin{proposition}[Falsador denso]
Si \(D\subsetneq X\) es denso, entonces
\[
  \pi_m(D)=X_m
  \qquad\text{para todo }m,
\]
de modo que \(D\) y \(X\) poseen las mismas sombras finitas.
\end{proposition}

\begin{proof}
Cada fibra cilíndrica no vacía \(\pi_m^{-1}(a)\) es abierta. Por densidad,
intersecta \(D\). Por tanto \(a\in\pi_m(D)\) para todo \(a\in X_m\).
\end{proof}

Las sombras finitas reconstruyen la clausura de un subconjunto, no el
subconjunto arbitrario. Éste es el punto exacto en que una identificación
\[
  \varprojlim_m\mathcal P(X_m)
  =
  \mathcal P(\varprojlim_mX_m)
\]
resultaría falsa.

\begin{figure}[H]
\centering
\begin{tikzpicture}[>=Latex,node distance=10mm and 17mm,
  box/.style={rounded corners,draw=hmtblue,fill=hmtlightblue,
    minimum width=38mm,minimum height=15mm,align=center},
  full/.style={rounded corners,draw=hmtviolet,fill=white,
    minimum width=40mm,minimum height=15mm,align=center}]
  \node[box] (fin) {\(\mathcal P(X_m)\)\\\scriptsize predicados finitos};
  \node[box,right=of fin] (clop) {\(\Clop(X)\)\\\scriptsize límite directo};
  \node[box,below=of clop] (closed) {\(\Closed(X)\)\\\scriptsize límite inverso};
  \node[full,right=of clop] (power) {\(\mathcal P(X)\)\\\scriptsize potencia plena};
  \draw[->,thick,hmtblue] (fin)--node[above]{\(p^{-1}\)}(clop);
  \draw[->,thick,hmtblue] (fin) to[bend right=24]
    node[below left]{\(p_*\)} (closed);
  \draw[->,thick,hmtgreen] (clop)--node[above]{\(\subset\)}(power);
  \draw[->,thick,hmtgreen] (closed)--node[below]{\(\subset\)}(power);
  \draw[dashed,very thick,hmtred] (closed.south east)
    .. controls +(8mm,-8mm) and +(-8mm,-8mm) ..
    node[midway,below=2mm]{\(\ne\)} (power.south west);
  \node[below=18mm of closed,text width=.82\textwidth,align=center] {
    \small Un denso propio y \(X\) tienen las mismas proyecciones finitas:
    la historia finita recupera clausura, no potencia plena.};
\end{tikzpicture}
\caption{Tres nociones de potencia que no deben aplanarse.}
\label{p3-83-c45-s2:fig:three-powers}
\end{figure}

\section{\(4\) exigencias de elección en el sistema proyectivo}

Sea \(p:X\twoheadrightarrow Y\). Conviene separar:

\begin{enumerate}
\item una \emph{sección conjuntista} \(s:Y\to X\) con \(ps=\id_Y\);
\item una \emph{familia natural} de secciones que conmuta con las
restricciones de una torre;
\item una \emph{sección algebraica} que además preserva las operaciones;
\item una \emph{sección efectiva} calculada por un algoritmo total.
\end{enumerate}

Toda sección natural, algebraica o efectiva es, olvidando estructura, una
sección conjuntista. Las tres propiedades adicionales son, en general,
incomparables: un algoritmo puede elegir representantes sin ser homomorfismo;
un homomorfismo puede no ser computable en la presentación elegida; una
familia natural puede no conservar la operación algebraica.

\begin{remark}[Condiciones de validez]
El axioma de elección garantiza una selección conjuntista dentro de su
dominio. No garantiza naturalidad, algebraicidad, efectividad ni coste
uniforme. La independencia de \(\mathsf{AC}\) respecto de \(\mathsf{ZF}\)
tampoco convierte estas cuatro exigencias en una sola: incluso en un modelo
con elección, la sección estructural necesita su propia prueba.
\end{remark}

\Needspace{36\baselineskip}
\section{Obstrucción de Pell a la periodicidad uniforme del sistema de prefijos}

El generador activo construye
\[
  K_\infty\cong C_4\times\mathbb Z_3
\]
y proyecciones finitas \(p_m:K_\infty\to K_m\). Puede elegirse un
representante de cada clase, pero no una sección homomorfa compatible con el
grupo.

\begin{proposition}
No existe homomorfismo \(s_m:K_m\to K_\infty\) con
\(p_ms_m=\id\) para el cociente que introduce torsión ternaria finita en
\(K_m\).
\end{proposition}

\begin{proof}
Una sección homomorfa enviaría un elemento de orden \(3^r\) de \(K_m\) a un
elemento del mismo orden en la componente \(\mathbb Z_3\). Pero el grupo
aditivo \(\mathbb Z_3\) es libre de torsión: si \(3^rx=0\), entonces \(x=0\).
La imagen no puede proyectar de nuevo sobre un elemento no nulo de orden
\(3^r\).
\end{proof}

La obstrucción no dice «no puede escogerse». Dice: no puede escogerse
preservando simultáneamente proyección y ley de grupo. La información del
carry que desaparece en el cociente reaparece como obstrucción a la sección.

\begin{remark}[Separación de resultados y alcance]
El continuo profinito y la obstrucción Pell son
\emph{resultados establecidos en las secciones precedentes}. Los dos teoremas de potencia y la tabla de
elecciones son \emph{construcciones explícitas del presente desarrollo}. El falsador denso impide
convertir la búsqueda incompleta de un subconjunto en una potencia ya
reconstruida.
\end{remark}
